\documentclass[../main]{subfiles}

\begin{document}

%% Give a brief description for the aims of the laboratory and
%% a summary of the results collected.
%% Please write this as constant prose in a way that
%% a fellow student, who does not know anything about the experiment,
%% can understand the approach.
\chapter{Introduction}
\label{ch:introduction}

\section{Four aims of the laboratory}

The aim of this laboratory is to build our own robotic car.
It is named as ``\textbf{Super Ultra Hyper Mega Advanced Senior Superior Elite Lobster}'',
short as ``\textbf{Lobster}''.
We have done \textbf{4} contributions to Lobster, which are recorded in
lab3-light seeking, lab5-line following, lab7-Obstacle avoiding and lab8-environment measurement.
Thus, Lobster could have the following functionalities:

\begin{enumerate}
\item Follow the light on the floor and turn its direction to wherever the light is.
\item Follow a 1 to 3 cm wide black line on the floor.
The line could be straight, bent or quarter.
\item Go straight and avoid obstacles.
When reaches obstacles such as walls, Lobster will turn to another direction to avoid crashing.
\item Measure temperature, humidity and atmospheric pressure,
print them and set a motor speed dependently.
\end{enumerate}

If Lobster has installed light seeking module, then it cannot track lines any more,
same as the other two functionalities.
In short, only one of these four modules works at a time.

\section{Results summary}

Basic approaches of the four experiment are completed.
We used two light-dependent resistors (LDR) to receive light change.
The car can recognize white light in the front within 15 cm. % TODO

\end{document}